\begin{proofbox}
\textbf{\textcolor{CProof}{四种证明要解决同一件事}}\par\smallskip
已知曲线上一切点 $T$，先证明过 $T$ 的切线有规定的形式；
再说明：若给定点 $P$ 在这条切线上，那么 $T$ 必位于 $P$ 的极线上。
这样，若有两个不同的切点，就得到切点弦。
\textbf{对于没有真实切点的点，极线仍是按式子给出的直线，并不是从切点弦推出来的。}

\medskip
\par
\begin{minipage}{\linewidth}
\textbf{\textcolor{CProof}{第一法｜门槛最低：圆的垂直关系 + 逐条代入检验}}\par\smallskip
\textbf{所用知识：}垂直、坐标、平方展开；不使用导数。

\begin{center}
\begin{tikzpicture}[x=1.3cm,y=1.3cm,>=Stealth,font=\small]
 \path[use as bounding box] (-1.5,-1.6) rectangle (8,2.25);
 \node at (.4,1.95) {单位圆：半径垂直切线};
 \draw[CStatement,very thick] (0,0) circle(1);
 \draw[orange!80!black,thick] (-.5,1.625)--(1.8,-.1);
 \draw[thick] (0,0)--(.6,.8);
 \draw[thick] (.51,.68)--(.63,.59)--(.72,.71);
 \fill (0,0) circle(1.5pt) node[below] {$O$};
 \fill (.6,.8) circle(1.5pt) node[above right] {$T'$};
 \fill (1.6,.05) circle(1.5pt) node[below right] {$Q'$};
 \draw[->,thick,gray] (2.1,.5)--(3.4,.5);
 \node[align=center] at (2.75,1.2) {$x=aX$\\$y=bY$};
 \begin{scope}[shift={(5.7,0)}]
  \node at (0,1.95) {椭圆：相切关系保留};
  \draw[CStatement,very thick] (0,0) ellipse (2 and 1);
  \draw[orange!80!black,thick] (-.8,1.55)--(2.1,.4625);
  \fill (1.2,.8) circle(1.5pt) node[above right] {$T(u,v)$};
  \fill (0,0) circle(1.5pt) node[below] {$O$};
 \end{scope}
 \node at (0,-1.4) {$X^2+Y^2=1$};
 \node at (5.7,-1.4) {$\dfrac{x^2}{a^2}+\dfrac{y^2}{b^2}=1$};
\end{tikzpicture}

\end{center}
左图 $T'=(u/a,v/b)$、$Q'=(x/a,y/b)$ 是伸缩前的点；
只有圆图使用半径垂直切线，椭圆图保留的是相切关系。

\end{minipage}
\par
\begin{description}[style=nextline,leftmargin=2.8em,labelsep=0.5em,itemsep=0.55em]
 \item[\textbf{步骤 1（从圆出发，看到“坐标相乘”）：}]
 圆 $x^2+y^2=R^2$ 的切点为 $T(u,v)$，则 $u^2+v^2=R^2$。
 切线上的任一点 $Q(x,y)$ 满足 $OT\perp TQ$，因此
 \[
 u(x-u)+v(y-v)=0
 \Longleftrightarrow ux+vy=R^2.
 \]
 \textbf{几何意义：}切线与半径垂直，使原先的平方项变成“切点坐标乘活动点坐标”。

 \item[\textbf{步骤 2（由圆到椭圆）：}]
 对椭圆 $x^2/a^2+y^2/b^2=1$，令 $X=x/a,\ Y=y/b$，
 则 $X^2+Y^2=1$。椭圆上切点 $T(u,v)$ 对应圆上点 $(u/a,v/b)$。
 圆的切线为
 \[
 \frac{u}{a}X+\frac{v}{b}Y=1.
 \]
 将 $X=x/a,Y=y/b$ 代回，得
 \[
 \boxed{\frac{ux}{a^2}+\frac{vy}{b^2}=1.}
 \]
 \textbf{看图核对：}伸缩保持直线和交点个数，所以圆的切线变为椭圆的切线；垂直关系一般不保留。

 \item[\textbf{步骤 3（双曲线：把候选切线代回去检查）：}]
 设 $T(u,v)$ 在 $x^2/a^2-y^2/b^2=1$ 上，即
 $u^2/a^2-v^2/b^2=1$。考虑过 $T$ 的直线
 \[
 \ell:\ \frac{ux}{a^2}-\frac{vy}{b^2}=1.
 \]
 它上面的点都可以写成
 \[
 x=u+a^2vt,\qquad y=v+b^2ut
 \quad(t\text{ 为实数}).
 \]
 因为 $t=0$ 给出 $T$，且把这组坐标代入 $\ell$ 左边得
 \[
 \frac{u(u+a^2vt)}{a^2}-\frac{v(v+b^2ut)}{b^2}
 =\frac{u^2}{a^2}-\frac{v^2}{b^2}+uvt-uvt=1,
 \]
 所以确实写出了 $\ell$ 上所有的点。再代入双曲线：
 \begin{align*}
 &\frac{(u+a^2vt)^2}{a^2}
       -\frac{(v+b^2ut)^2}{b^2}-1\\
 &=(\tfrac{u^2}{a^2}-\tfrac{v^2}{b^2}-1)
    +(2uvt-2uvt)+(a^2v^2-b^2u^2)t^2\\
 &=-a^2b^2t^2.
 \end{align*}
 最后一步用到了
 $b^2u^2-a^2v^2=a^2b^2$。
 由于 $a,b>0$，只有 $t=0$ 时才落在双曲线上，所以 $\ell$ 与双曲线只在 $T$ 接触，确为切线。
 \textbf{几何意义：}沿这条直线离开 $T$，曲线方程左边减去右边的值始终是负数；接触只发生在 $T$。

 \item[\textbf{步骤 4（抛物线：同样只用平方展开）：}]
 设 $T(u,v)$ 在 $y^2=2px$ 上，故 $v^2=2pu$。
 考虑过 $T$ 的直线 $vy=p(x+u)$。
 它上面的点都可以写成
 \[
 x=u+vt,\qquad y=v+pt.
 \]
 一方面
 $v(v+pt)=v^2+pvt=p(2u+vt)=p((u+vt)+u)$，
 说明这些点在候选直线上；另一方面
 \begin{align*}
 y^2-2px
  &=(v+pt)^2-2p(u+vt)\\
  &=v^2+2vpt+p^2t^2-2pu-2pvt\\
  &=p^2t^2.
 \end{align*}
 因为 $p>0$，只有 $t=0$ 时同时在抛物线上，故
 \[
 \boxed{vy=p(x+u)}
 \]
 是切线。\textbf{特别地，顶点 $T(0,0)$ 的切线是 $x=0$，并没有被斜率形式漏掉。}

 \item[\textbf{步骤 5（把切线转成切点弦）：}]
 以椭圆为例，若从 $P(x_0,y_0)$ 作两条切线，切点为
 $T_1(u_1,v_1)$、$T_2(u_2,v_2)$，
 因为 $P$ 在这两条切线上，所以
 \[
 \frac{u_1x_0}{a^2}+\frac{v_1y_0}{b^2}=1,
 \qquad
 \frac{u_2x_0}{a^2}+\frac{v_2y_0}{b^2}=1.
 \]
 把 $(u_1,v_1),(u_2,v_2)$ 看成待求点 $(x,y)$，便知两个切点都位于
 \[
 \boxed{\frac{x_0x}{a^2}+\frac{y_0y}{b^2}=1.}
 \]
 两个\textbf{不同的点}确定一条直线，因此此线就是切点弦。
 双曲线和抛物线只需把本法步骤 3、4 的切线公式作同样的坐标交换即可。
\end{description}

\medskip\hrule\medskip
\par
\begin{minipage}{\linewidth}
\textbf{\textcolor{CProof}{第二法｜速度最快：已知切线公式，直接交换两组坐标}}\par\smallskip
\textbf{适用场景：}考试中已经知道曲线上一点的切线方程，题目要从外点求切点弦。

\begin{center}
\begin{tikzpicture}[x=1.65cm,y=1.65cm,>=Stealth,font=\small]
 \path[use as bounding box] (-2.4,-1.65) rectangle (4.7,1.65);
 \draw[CStatement,very thick] (0,0) ellipse (2 and 1);
 \draw[orange!80!black,thick] (.3,1.0681)--(4,0)--(.3,-1.0681);
 \draw[teal!80!black,very thick] (1,-1.3)--(1,1.3);
 \fill (4,0) circle(1.8pt) node[right] {$P(x_0,y_0)$};
 \fill (1,.8660254) circle(1.8pt) node[above right] {$T_1(u_1,v_1)$};
 \fill (1,-.8660254) circle(1.8pt) node[below right] {$T_2(u_2,v_2)$};
 \node[CStatement] at (-.8,1.45) {椭圆};
 \node[teal!80!black,align=center] at (-.7,-1.4) {两切点满足同一条极线方程};
\end{tikzpicture}

\end{center}
观察橙色两切线都过 $P$：分别代入 $P$，便得到两个切点满足的同一个方程。
\end{minipage}
\par
\begin{description}[style=nextline,leftmargin=2.8em,labelsep=0.5em,itemsep=0.45em]
 \item[\textbf{步骤 1（切点是未知的）：}]
 设切点 $T(u,v)$，外点 $P(x_0,y_0)$。标准椭圆在 $T$ 处的切线为
 $ux/a^2+vy/b^2=1$。
 \item[\textbf{步骤 2（因为 $P$ 在切线上）：}]
 立即写出
 \[
 \frac{ux_0}{a^2}+\frac{vy_0}{b^2}=1.
 \]
 此式没有求出 $u,v$，但已给出\textbf{所有可能切点}必须满足的条件。
 \item[\textbf{步骤 3（把 $u,v$ 改写成 $x,y$）：}]
 \[
 \boxed{\frac{x_0x}{a^2}+\frac{y_0y}{b^2}=1}.
 \]
 双曲线直接把 $+$ 换成 $-$；抛物线由切线 $vy=p(x+u)$
 和 $P$ 在切线上得到 $v y_0=p(x_0+u)$，即
 \[
 \boxed{y_0y=p(x+x_0)}.
 \]
 对一般二次曲线，同样使用第四法证明的切线公式，
 交换 $P,T$ 的坐标后就得到结论三的等式。
\end{description}
\textbf{速度来自哪里：}只写一次“切线经过 $P$”，不用先联立求两个切点。
\textbf{限制：}这种做法证明的是\textbf{有真实切点时}的切点弦，不能据此断言任何点都有两条切线。

\medskip\hrule\medskip
\par
\begin{minipage}{\linewidth}
\textbf{\textcolor{CProof}{第三法｜步骤最少：把相切看作二次方程出现重根}}\par\smallskip
\textbf{特点：}只用一个一元二次式，同时适用于含 $xy$ 项的一般曲线。

\begin{center}
\begin{tikzpicture}[x=1.6cm,y=1.6cm,>=Stealth,font=\small]
 \path[use as bounding box] (-1.7,-1.4) rectangle (6,2.2);
 \draw[CStatement,thick] (0,0) circle(1);
 \draw[orange!80!black,thick,->] (-.9,1)--(2.5,1);
 \fill (2,1) circle(1.8pt) node[above] {$P$: $t=0$};
 \fill (0,1) circle(1.8pt) node[above] {$T$: $t=1$};
 \node at (.4,-1.3) {$Q(t)=P+t(T-P)$};
 \begin{scope}[shift={(4,0)}]
  \draw[->,gray] (-1,0)--(1.75,0) node[right] {$t$};
  \draw[->,gray] (-.8,-.2)--(-.8,1.65);
  \node at (.85,1.25) {$f(t)/\alpha$};
  \draw[CProof,very thick,domain=-.25:1.75,samples=60] plot (\x,{(\x-1)^2});
  \fill (1,0) circle(1.8pt) node[below] {$1$（重根）};
  \node at (.35,2.1) {$f(t)=\alpha(t-1)^2$};
 \end{scope}
\end{tikzpicture}

\end{center}
左图用圆 $G=x^2+y^2-1$ 演示：$P=(2,1)$、$T=(0,1)$，
故 $Q(t)=(2-2t,1)$，$f(t)=4(t-1)^2$。右图画 $f(t)/\alpha$，
让“切点对应 $t=1$ 的重根”直接可见。
\end{minipage}
\par
\begin{description}[style=nextline,leftmargin=2.8em,labelsep=0.5em,itemsep=0.5em]
 \item[\textbf{步骤 1（用 $t$ 表示 $P$ 到切点 $T$ 的直线）：}]
 记
 \[
 G(x,y)=Ax^2+Bxy+Cy^2+Dx+Ey+F,
 \quad T(u,v),\ P(x_0,y_0).
 \]
 当 $P\ne T$，直线 $PT$ 上一点写成
 \[
 (x,y)=(x_0+t(u-x_0),\ y_0+t(v-y_0)).
 \]
 代入曲线方程，可整理为
 \[
 f(t)=\alpha t^2+\beta t+\gamma,
 \]
 其中
 \[
 \begin{aligned}
 \gamma={}&G(x_0,y_0),\\
 \beta={}&(2Ax_0+By_0+D)(u-x_0)\\
             &+(Bx_0+2Cy_0+E)(v-y_0).
 \end{aligned}
 \]
 \item[\textbf{步骤 2（切点就是重根 $t=1$）：}]
 因为 $PT$ 在 $T$ 处相切，$t=1$ 为二次方程的重根，故
 \[
 f(t)=\alpha(t-1)^2=\alpha t^2-2\alpha t+\alpha.
 \]
 比较系数：$\beta=-2\alpha,\ \gamma=\alpha$，从而
 \[
 \boxed{\gamma+\frac{\beta}{2}=0}.
 \]
 \item[\textbf{步骤 3（把该式整理为极线）：}]
 将第一步中的 $\gamma,\beta$ 代入左边：
 \begin{align*}
 \gamma+\frac{\beta}{2}
 ={}&Ax_0^2+Bx_0y_0+Cy_0^2+Dx_0+Ey_0+F\\
 &+\left(Ax_0+\frac{By_0}{2}+\frac D2\right)(u-x_0)\\
 &+\left(\frac{Bx_0}{2}+Cy_0+\frac E2\right)(v-y_0)\\
 ={}&Ax_0u+\frac B2(x_0v+y_0u)+Cy_0v\\
 &+\frac D2(x_0+u)+\frac E2(y_0+v)+F.
 \end{align*}
 因此切点 $T(u,v)$ 满足
 \[
 Ax_0u+\frac B2(x_0v+y_0u)+Cy_0v
 +\frac D2(x_0+u)+\frac E2(y_0+v)+F=0.
 \]
 令 $(u,v)$ 为待求切点 $(x,y)$，即为极线方程。
\end{description}
\textbf{边界说明：}如果 $P=T$，不需作参数代入：$P$ 位于曲线上，
 所以把 $(x,y)=(x_0,y_0)$ 代入极线等式，就还原为 $G(P)=0$；
 该直线确为切线，由第四法直接证明。

\medskip\hrule\medskip
\par
\begin{minipage}{\linewidth}
\textbf{\textcolor{CProof}{第四法｜常规通用：设过切点的直线，联立并比较重根}}\par\smallskip
\textbf{优点：}从头推导一般二次曲线的切线公式，不必事先记住任何一种圆锥曲线的切线。

\begin{center}
% Circle is a concrete instance of a general conic; T=(0,1).
\begin{tikzpicture}[x=1.6cm,y=1.6cm,>=Stealth,font=\small]
 \path[use as bounding box] (-2,-1.55) rectangle (5.3,2.2);
 \draw[CStatement,very thick] (0,0) circle(1);
 \draw[orange!80!black,very thick] (-1.65,1)--(2.3,1);
 \fill (0,1) circle(1.8pt) node[above left] {$T(u,v)$};
 \draw[->,thick,CProof] (.35,1)--(1.7,1);
 \node[CProof] at (.95,1.4) {方向 $(r,s)$};
 \draw[->,thick,teal!80!black] (0,1)--(0,2) node[right] {$\boldsymbol n$};
 \draw (0,1.15)--(.15,1.15)--(.15,1);
 \node at (3.15,1.6) {$\boldsymbol n=(2Au+Bv+D,$};
 \node at (3.15,1.15) {$Bu+2Cv+E)$};
 \node[align=center] at (3.15,.25) {切线方向与 $\boldsymbol n$ 垂直\\一次项系数为 $0$};
 \node at (1.4,-1.35) {以圆为例：$T$ 对应 $t=0$，沿切线移动};
\end{tikzpicture}

\end{center}
图中用圆上的切点演示一般关系；$\boldsymbol n$ 是垂直于切线的向量。
沿方向 $(r,s)$ 移动时，一次项消失，对应切点 $t=0$ 为重根。
\end{minipage}
\par
\begin{description}[style=nextline,leftmargin=2.8em,labelsep=0.5em,itemsep=0.55em]
 \item[\textbf{步骤 1（设定切点和直线方向）：}]
 在 $\Gamma:G(x,y)=0$ 上取点 $T(u,v)$，则
 \[
 Au^2+Buv+Cv^2+Du+Ev+F=0.
 \]
 过 $T$ 的一条直线可表示为
 \[
 x=u+rt,\qquad y=v+st,
 \quad (r,s)\ne(0,0).
 \]
 这里 $(r,s)$ 只表示沿直线向前走的方向，\textbf{没有用斜率}，因此竖直线也包含在内。

 \item[\textbf{步骤 2（代入曲线，逐项展开）：}]
 \begin{align*}
 0={}&A(u+rt)^2+B(u+rt)(v+st)\\
    &+C(v+st)^2+D(u+rt)+E(v+st)+F\\
  ={}&\underbrace{(Au^2+Buv+Cv^2+Du+Ev+F)}_{=0}\\
   &+t\bigl[(2Au+Bv+D)r+(Bu+2Cv+E)s\bigr]\\
   &+t^2(Ar^2+Brs+Cs^2).
 \end{align*}
 \textbf{几何意义：}上式就是“直线与曲线的交点”所满足的一元二次方程。
 切点对应 $t=0$；相切时两个交点重合为 $t=0$，所以一次项必须消失。

 \item[\textbf{步骤 3（让 $t=0$ 成为重根）：}]
 因为一次项系数为零，故切线的方向 $(r,s)$ 满足
 \[
 \boxed{(2Au+Bv+D)r+(Bu+2Cv+E)s=0}.
 \]
 过切点 $T$ 的任一点 $Q(x,y)$ 与 $T$ 的坐标差为
 $(x-u,y-v)$，也沿着这条切线，于是可以把
 $r,s$ 分别换成 $x-u,y-v$，得到
 \[
 (2Au+Bv+D)(x-u)+(Bu+2Cv+E)(y-v)=0.
 \]
 \textbf{几何意义：}系数对 $(2Au+Bv+D,\ Bu+2Cv+E)$
 与切线方向垂直，就是这条切线的一个法向量。

 \item[\textbf{步骤 4（展开并利用切点在曲线上）：}]
 上式展开为
 \[
 (2Au+Bv+D)x+(Bu+2Cv+E)y
 -(2Au^2+2Buv+2Cv^2+Du+Ev)=0.
 \]
 由 $G(u,v)=0$ 得
 \[
 Au^2+Buv+Cv^2=-Du-Ev-F,
 \]
 所以
 \[
 2Au^2+2Buv+2Cv^2+Du+Ev=-Du-Ev-2F.
 \]
 代回并把各项除以 $2$：
 \[
 \boxed{\begin{aligned}
 &Aux+\frac B2(uy+vx)+Cvy\\
 &\qquad+\frac D2(x+u)+\frac E2(y+v)+F=0.
 \end{aligned}}
 \]
 这就是一般二次曲线在 $T$ 处的切线方程。

 \item[\textbf{步骤 5（由切线推出 $P$ 的极线）：}]
 若 $P(x_0,y_0)$ 在上述切线上，就有
 \[
 Ax_0u+\frac B2(uy_0+vx_0)+Cy_0v
 +\frac D2(x_0+u)+\frac E2(y_0+v)+F=0.
 \]
 它与把 $T=(u,v)$ 代入 $P$ 的极线公式得到的式子完全相同。
 因此所有从 $P$ 引出的真实切线的切点都在 $P$ 的极线上；
 \textbf{若有两个不同切点，极线就是过这两个点的直线}。
\end{description}

\medskip
\textbf{\textcolor{CProof}{最后统一检查}}\par\smallskip
\begin{enumerate}[leftmargin=*,itemsep=0.3em]
 \item $P$ 的极线为普通直线，当且仅当
 $2Ax_0+By_0+D$ 与 $Bx_0+2Cy_0+E$ 不同时为零。
 \item 若 $P$ 在曲线上，以上第四法证明：极线就是 $P$ 处切线。
 \item 若 $P$ 的极线与曲线有两个\textbf{不同的实交点}，第五步反过来说明 $P$ 落在这两个交点各自的切线上，所以它们正是切点。
 \item 若极线没有两个不同的实交点，就不能称它为实切点弦；对于根本不能形成直线的特殊点，更不能套用“切点弦”结论。
\end{enumerate}
\end{proofbox}
